\begin{center}\begin{tikzpicture}[font=\small,>=stealth,line join=round]
\draw[gray!40,rounded corners=3pt] (0,0) rectangle (14.2,8.3);\node at (7.1,7.700000000000001) {例6：外层只需接受内层实际给出的数};
\begin{scope}[shift={(4.1,2.2716216216216214)},x=1.0714285714285714cm,y=1.243243243243243cm]\draw[->,gray] (-2.8,0)--(2.8,0) node[right]{$x$};\draw[->,gray] (0,-0.5)--(0,3.2) node[above]{$y$};\draw[blue,thick,solid,domain=-2.5:2.5,samples=180] plot(\x,{sqrt((\x)^2+1)});\fill[orange!85!black] (0,1) circle(2.2pt);\end{scope}\node[anchor=west,align=left,text width=5.0cm] at (8.4,4.0) {$F=\sqrt{x^2+1}$\\[10pt]内层实际输出 $[1,+\infty)$\\[10pt]全部能开平方\\[10pt]整体在 $\mathbb R$ 上为偶函数};
\node[text width=12.6cm,align=center] at (7.1,.65) {外层完整定义域不必左右对称；本题只使用非负的那部分输入。};
\end{tikzpicture}\end{center}
